\unit{Driving, cooling and trapping atoms with lasers}
\unitref{Week 10; 2025 Q2; Zeeman slower tutorial}
\subsection{A laser couples two atomic levels}
Near an isolated optical transition,
model the atom with \(\ket g,\ket e\).
A classical electric field
\(\mathbf E(t)=\mathbf E_0\cos(\omega_Lt)\)
couples through \(-\mathbf d\cdot\mathbf E(t)\).
The \emph{Rabi frequency} measures
interaction strength:
\[
\boxed{\Omega=\frac{|\bra e\mathbf d\cdot
\mathbf E_0\ket g|}{\hbar}}
\]
with factors of 2 depending on whether
\(E_0\) denotes the full real-field
amplitude or its positive-frequency
part. State the convention in a
pulse-timing problem.
In a rotating frame and under
the rotating-wave approximation,
a convenient effective Hamiltonian is
\[
H_{\rm rot}=\frac{\hbar}2
\begin{pmatrix}-\Delta&\Omega\\
\Omega&\Delta\end{pmatrix},
\quad
\Delta=\omega_L-\omega_{eg}.
\]
The generalised Rabi frequency is
\(\Omega_{\rm eff}=\sqrt{\Omega^2+\Delta^2}\).

\subsection{Pulse area, pi and pi-over-two pulses}
For a resonant pulse with constant
coupling phase and time-dependent
amplitude \(\Omega(t)\),
the pulse area is
\(\Theta=\int\Omega(t)dt\).
An atom starting in \(\ket g\) has
\[
\boxed{P_e=\sin^2(\Theta/2).}
\]
A \(\pi\)-pulse transfers all population
to the excited state in the ideal
lossless model; a \(\pi/2\)-pulse
creates an equal-amplitude superposition.
Detuning prevents perfect inversion for
a constant-frequency rectangular pulse,
because the largest excited probability is
\(\Omega^2/(\Omega^2+\Delta^2)\).

\subsection{Why light can exert a force}
A photon carries momentum \(\hbar\mathbf k\).
Absorption transfers this directed momentum
to an atom. Spontaneous photons are emitted
in different directions, so their average
momentum along a fixed axis may vanish,
though their variance gives recoil heating.
For a driven two-level atom with excited-state
decay rate \(\Gamma\), a common steady-state
scattering rate is
\[
\boxed{R_{\rm sc}=\frac{\Gamma}{2}
\frac{s}{1+s+(2\Delta/\Gamma)^2}},
\qquad s=I/I_{\rm sat}.
\]
At high saturation on resonance,
\(R_{\rm sc}\to\Gamma/2\), giving
\[
\boxed{F_{\max}=\hbar k\Gamma/2,\qquad
a_{\max}=F_{\max}/M.}
\]
There is a trade-off: fast-scattering
broad transitions give strong slowing
forces but generally higher Doppler
cooling limits than narrow lines.

\subsection{Optical molasses and the Doppler shift}
An atom moving toward one of two
counter-propagating red-detuned beams
sees that beam Doppler-shifted
closer to resonance. It absorbs
more photons from the direction
opposing its motion, creating a
damping force near rest,
\(F\simeq-\beta v\).
In a simple two-level Doppler
cooling model, the minimum temperature is
\[
\boxed{T_D=\frac{\hbar\Gamma}{2k_B}
=\frac{\hbar}{2k_B\tau}},
\]
where \(\tau=1/\Gamma\).
The temperature comes from balancing
velocity damping against random
photon-recoil diffusion; the kinetic
energy never falls to exactly zero.

\subsection{A Zeeman slower follows changing velocity}
A fast atomic beam has a large Doppler
shift, which changes as atoms slow.
A position-dependent magnetic field
Zeeman-shifts the atomic transition
to maintain resonance with the
counter-propagating laser.
If the designed deceleration is constant,
\[
v(z)=\sqrt{v_0^2-2az},
\quad L=\frac{v_0^2}{2a},
\]
and, for a suitable sign convention,
the resonance field has the shape
\[
\boxed{B(z)=B_{\rm offset}
+B_0\sqrt{1-z/L}.}
\]
This is the origin of the characteristic
square-root slower profile, not merely
a curve fitted to experimental data.
An arbitrary field that reaches zero
need not slow atoms successfully:
the corresponding resonant velocity
must be dynamically followable.
Specifically,
\[
\boxed{a_{\rm req}(z)=
-v_{\rm res}(z)\frac{dv_{\rm res}}{dz}}
\]
must stay within the available scattering
acceleration and capture linewidth.

\begin{worked}{Estimate ytterbium stopping distance}
For the 2025 exam data,
\(M=1.73\times10^{-25}\,\mathrm{kg}\),
\(T=300\,\mathrm{K}\),
\(\lambda=399\,\mathrm{nm}\),
\(\tau=5.5\,\mathrm{ns}\),
choose the 3D rms speed
\(v_0=\sqrt{3k_BT/M}\approx268\,\mathrm{m/s}\).
At half the maximum scattering force,
\[
a=\frac{h}{4M\lambda\tau}
\approx4.36\times10^5\,\mathrm{m/s^2}.
\]
Hence
\[
L=\frac{v_0^2}{2a}\approx8.23\,\mathrm{cm},
\qquad
N_\gamma=\frac{Mv_0\lambda}{h}
\approx2.79\times10^4.
\]
The Doppler limit from this linewidth
is \(T_D\approx0.694\,\mathrm{mK}\).
The thermal speed convention affects
the stopping distance; state it explicitly.
\end{worked}

\subsection{Magneto-optical traps and selection rules}
A magneto-optical trap (MOT) uses
opposing red-detuned laser beams
with suitable circular polarisations
and a magnetic-field gradient.
Position-dependent Zeeman shifts
make one opposing beam more resonant,
producing a restoring force, while the
Doppler effect supplies damping.
Polarisation selection rules connect
light to atomic magnetic substates:
\(\pi\)-polarised light has \(\Delta m=0\);
\(\sigma^\pm\) components have
\(\Delta m=\pm1\), for a consistent
quantisation axis and photon helicity.
Optical pumping can redistribute
population among those substates.

\begin{checkpoint}
Rabi oscillations describe coherent drive.
Spontaneous scattering describes dissipative
momentum kicks. Doppler molasses damps
velocity; a MOT also restores position.
Zeeman slowing maintains resonance
while the beam loses speed.
\end{checkpoint}
